% ==========================================================
\section{ANTLR Parser Integration and Metamodel Extensions for BP}
\label{sec:impl_uvl_metamodel}
% ==========================================================

This section describes the concrete realization of the conceptual \gls{bp} integration we introduce in Section~\ref{sec:concept_bp_integration}.
It covers the grammar extensions, parser adaptations, and metamodel classes that constitute the implementation.

% ----------------------------------------------------------
\subsection{ANTLR Grammar Extensions}
\label{subsec:impl_uvl_metamodel_antlr}
% ----------------------------------------------------------

The grammar modifications are confined to \texttt{UVLParser.g4} in the \gls{antlr} grammar module of the \gls{uvl} metamodel component.
They realize the opt-in \gls{bp} extension described at the concept level while leaving all remaining productions unchanged, thereby preserving backward compatibility.
The extensions comprise three targeted additions, shown in Listing~\ref{lst:antlr-extension}.
First, we augment the \texttt{features} rule to allow up to two additional top-level features, \texttt{Env} and \texttt{Config}.
Second, we extend the \texttt{constraint} rule with a new labeled alternative \texttt{bpConstraint}, which groups all \gls{bp}-specific constraint types under a single production to maintain an organized grammar structure.
Third, within \texttt{bpConstraint}, a dedicated \texttt{conflictingConstraint} production handles the multi-reference case of mutually conflicting b-events.
To ensure deterministic lexing, we introduce a dedicated keyword token for each \texttt{bp} construct, as shown in Listing~\ref{lst:bp-lexer-extension}.

\lstinputlisting[
    float,
    floatplacement=ht,
    caption={ANTLR grammar extensions for BP integration.},
    label={lst:antlr-extension}
]{code/metamodel/uvl_bp_grammar_snippet.g4}

\lstinputlisting[
    float,
    floatplacement=ht,
    caption={ANTLR lexer extensions for BP integration.},
    label={lst:bp-lexer-extension}
]{code/metamodel/bp_lexer_extensions.g4}

% ----------------------------------------------------------
\subsection{Parser and Factory Adaptations}
\label{subsec:impl_uvl_metamodel_parser}
% ----------------------------------------------------------

We adapt the parsing logic in two places: the \gls{antlr} listener and the model factory.
The \texttt{ModelType} enumeration controls model variant selection, introduced to support distinct parsing strategies under a single parser entry point, and may be extended to accommodate additional model variants in the future.
When \texttt{ModelType.BP} is active, the \texttt{UVLListener} constructor instantiates a \texttt{BPFeatureModel}.
Otherwise, it falls back to the standard \texttt{FeatureModel}, as shown in Listing~\ref{lst:uvl-listener-constructor}.

\lstinputlisting[
    language=Java,
    float,
    floatplacement=ht,
    caption={\texttt{UVLListener} constructor with model type selection.},
    label={lst:uvl-listener-constructor}
]{code/metamodel/UVLListener_constructor.java}

\gls{antlr} generates paired \texttt{enterX} and \texttt{exitX} callbacks for each grammar rule.
A so-called \texttt{ParseTreeWalker} invokes \texttt{enterX} before visiting a node's children and \texttt{exitX} after all child callbacks have completed.
The \texttt{enter} methods initialize objects and push them onto a stack, while \texttt{exit} methods finalize those objects using data populated by child callbacks~\cite{parr_definitive_2013}.

The additional top-level feature rule illustrates this pattern concretely.
Its grammar production is \texttt{reference attributes? NEWLINE}, meaning that when the walker enters the rule, the \texttt{reference} and \texttt{attributes} subtrees are accessible but not yet resolved.
Both listener methods begin with a call to \texttt{rejectIfBaseModel}, which enforces the opt-in nature of \gls{bp} parsing.
When the active model type is not \texttt{ModelType.BP}, any \gls{bp}-specific construct encountered in the listener is rejected with a parse error, and the method returns immediately.

Accordingly, \texttt{enterEnvConfigFeature} creates a \texttt{Feature} and pushes it onto the \texttt{featureStack}.
Child listeners populate the feature instance from its child parse nodes, namely the \texttt{reference} and optional \texttt{attributes} subtrees.
When we invoke \texttt{exitEnvConfigFeature}, the \texttt{Feature} is popped from the stack and assigned to either the \texttt{Env} or \texttt{Config} field.
Any feature name other than these two gets rejected.
Listing~\ref{lst:bp-top-level-listener} shows the corresponding listener methods.

\lstinputlisting[
    language=Java,
    float,
    floatplacement=ht,
    caption={Listener methods for BP top-level features (Env / Config).},
    label={lst:bp-top-level-listener}
]{code/metamodel/BP_top_level_listener.java}

\gls{bp} constraints follow a similar enter/exit pattern.
The enter method instantiates the appropriate constraint object, and the exit method finalizes it and registers it with the model.
For \texttt{ConflictingConstraint}, which carries multiple references, the listener accumulates all reference objects into a list during child callbacks and passes that list to the constructor upon exit.
The full implementations are available in the accompanying source code~\cite{ruider_xxnicksdaxxma-nick-ruider_2026}.

% ----------------------------------------------------------
\subsection{Metamodel Extensions}
\label{subsec:impl_uvl_metamodel_java_extension}
% ----------------------------------------------------------

The metamodel classes introduced in this subsection are the same ones instantiated by the listener described in the previous Section~\ref{subsec:impl_uvl_metamodel_parser}.
They provide the structural definitions that the listener populates during parsing.

The central addition is \texttt{BPFeatureModel}, which extends \texttt{FeatureModel} and introduces two dedicated fields for the \texttt{Env} and \texttt{Config} top-level features.
Both are represented as standard \texttt{Feature} objects and are distinguished semantically by their type attribute definition.
Listing~\ref{lst:bp-feature-model} shows an excerpt of the class structure.

\lstinputlisting[
    language=Java,
    float,
    floatplacement=ht,
    caption={\texttt{BPFeatureModel} class with \texttt{Env} and \texttt{Config} top-level features.},
    label={lst:bp-feature-model}
]{code/metamodel/BPFeatureModel.java}

\lstinputlisting[
    language=Java,
    float,
    floatplacement=ht,
    caption={\texttt{ConflictingConstraint} class for multi-reference BP constraints (excerpt).},
    label={lst:conflicting-constraint}
]{code/metamodel/ConflictingConstraint.java}

Another example of a metamodel addition is the new \texttt{ConflictingConstraint}.
Rather than holding a single reference, it manages an ordered list of \texttt{VariableReference} objects.
The class is shown in Listing~\ref{lst:conflicting-constraint}.

All remaining \gls{bp}-specific constraint additions follow the same structural pattern: a grammar production, a listener that constructs the constraint object, and registration with the model upon exit.
Single-reference \gls{bp} constraints extend \texttt{AbstractBPEventConstraint}, which is a shared base class that encapsulates the common single-event reference structure.
Further implementation details are available in the source code~\cite{ruider_xxnicksdaxxma-nick-ruider_2026}.

In summary, we realize the \gls{bp} integration through targeted extensions at each layer.
All changes are strictly additive with respect to the standard \gls{uvl} pipeline, yielding a \gls{bp}-capable parsing infrastructure that is backward-compatible with existing \gls{uvl} models and structurally aligned with the conceptual architecture developed in this thesis.